\section{为什么要单独讲硬件}

操作系统教材常把硬件压缩成几张框图：CPU、内存、总线、磁盘、网卡。这样会让后面的概念失去来源。系统调用为什么必须切栈，缺页为什么可以重试，同步为什么要原子指令，DMA 为什么要屏障，调度为什么关心 cache 和 TLB，这些问题都要回到硬件怎样由电路一步一步组成。

本章只沿 x86-64 主线解释通用计算机。目标不是把芯片设计课完整搬进来，而是把 OS 必须依赖的硬件机制讲到可推理、可画路径、可分析时延、可解释失败。每个硬件模块都按六个问题拆：

\begin{longtable}{p{0.18\linewidth}p{0.74\linewidth}}
\toprule
问题 & 需要回答什么 \\
\midrule
输入 & 这个模块从哪里拿 bit、地址、控制信号或时钟 \\
输出 & 它把什么结果交给下一级，结果什么时候有效 \\
组合路径 & 输出是否只由当前输入决定，最长逻辑路径是什么 \\
状态 & 它是否保存历史，状态在什么时钟边沿更新 \\
时延 & 它的关键路径如何限制频率、吞吐或响应延迟 \\
OS 关系 & 它如何影响异常、特权、调度、内存、I/O 或安全 \\
\bottomrule
\end{longtable}

有了这张拆解表，硬件就不再是一堆名词。加法器不是“能加数”的盒子，而是一串产生 sum/carry 的门；寄存器不是变量，而是时钟边沿采样的状态；ALU 不是魔法函数，而是多个组合电路加一个选择器；CPU 不是一个大循环，而是寄存器、组合逻辑、控制信号、时钟和异常入口共同组成的状态机。

\section{晶体管到逻辑门：先得到可靠的 0 和 1}

数字电路的第一层抽象是把连续电压分成两个稳定区域：低电平表示 0，高电平表示 1。晶体管可以看成受输入控制的开关。很多开关组合起来，就能实现逻辑门。真实芯片要考虑噪声容限、功耗、扇出、电容负载和工艺波动；教材模型先抓住一件事：输入稳定后，输出经过传播延迟也会稳定。

\begin{longtable}{p{0.16\linewidth}p{0.30\linewidth}p{0.42\linewidth}}
\toprule
门 & 布尔函数 & 为什么重要 \\
\midrule
NOT & 输出输入的反值 & 构造反相、使能取反、条件翻转 \\
NAND & 不是两个输入都为 1 & 功能完备，可构造任意组合逻辑 \\
NOR & 不是任一输入为 1 & 功能完备，适合某些控制逻辑 \\
AND & 两个条件都为 1 & 权限位、valid 位、enable 位同时满足 \\
OR & 至少一个条件为 1 & 多个异常源、多个 pending 位汇总 \\
XOR & 两个输入不同 & 加法器 sum、奇偶校验、比较辅助 \\
\bottomrule
\end{longtable}

“功能完备”的意思是，只用 NAND 或只用 NOR 也能搭出所有布尔函数。例如 NOT 可以用 NAND 的两个输入接同一个信号得到；AND 可以先 NAND 再 NOT；OR 可以用德摩根律得到。芯片设计不会真的每处都只手写 NAND，但理解 NAND/NOR 能说明复杂硬件最终还是由可组合的简单门构成。

\begin{lstlisting}
not_a = nand(a, a)
a_and_b = not(nand(a, b))
a_or_b = nand(not(a), not(b))
\end{lstlisting}

门有传播延迟。若一个 NAND 延迟记为 \code{t\_nand}，三个门串联的路径至少要等三个门都稳定。组合逻辑的输出不能在输入改变后立即采样，必须给信号传播时间。这个事实会一路推导出时钟周期、关键路径、流水线和 CPU 频率上限。

\section{组合逻辑：没有记忆的计算网络}

组合逻辑只由当前输入决定输出，不保存历史。加法器、比较器、译码器、选择器、优先级编码器都属于组合逻辑。组合逻辑适合表达“现在这些 bit 应该产生什么结果”，但不能表达“上一周期执行到哪里”。

\subsection{译码器和编码器}

译码器把一个二进制编号变成多根 one-hot 信号。例如 3 位输入可以选择 8 个输出中的一个。寄存器文件要根据寄存器编号选择某个寄存器，指令译码器要根据 opcode 选择控制路径，中断控制器要根据 vector 找入口。

\begin{lstlisting}
decode3(x):
  out[0] = (x == 0)
  out[1] = (x == 1)
  ...
  out[7] = (x == 7)
\end{lstlisting}

编码器反过来，把多根输入压成编号。优先级编码器还要规定多个输入同时为 1 时选谁。例如多个中断 pending 时，中断控制器要根据优先级选择一个 vector 投递给 CPU。

\begin{longtable}{p{0.20\linewidth}p{0.34\linewidth}p{0.34\linewidth}}
\toprule
模块 & 输入输出 & OS 相关影子 \\
\midrule
寄存器编号译码 & register id 到读写选择线 & syscall ABI 规定参数在哪些寄存器 \\
opcode 译码 & 指令字节到控制信号 & \code{syscall}、特权指令和非法指令检查 \\
vector 选择 & pending 中断到 vector & timer、IPI、device IRQ 分发 \\
页表索引拆分 & 虚拟地址位段到 PML4/PDPT/PD/PT index & page fault handler 能定位 PTE \\
\bottomrule
\end{longtable}

\subsection{多路选择器 MUX}

MUX 根据选择信号从多个输入中选一个输出。CPU 里到处都是 MUX：ALU 的第二个操作数来自寄存器还是立即数，写回值来自 ALU 还是 load 数据，下一 RIP 来自顺序地址、分支目标、异常入口还是返回地址。

\begin{lstlisting}
writeback_value = mux(
  from_alu: alu_result,
  from_memory: load_data,
  from_immediate: sign_extended_imm,
  from_status: flags_value
)
\end{lstlisting}

MUX 的代价是选择路径也有延迟。输入很多时，MUX 会变深或变宽，时序更难。x86-64 的复杂寻址模式、复杂写回来源和丰富控制状态，都会增加前端译码和后端选择网络的复杂度。OS 不直接控制 MUX，但系统调用、异常、返回用户态都依赖“下一 RIP 只能由受控硬件路径选择”。

\section{加法器：从一位到 64 位}

\subsection{半加器和全加器}

一位加法先从真值表开始。两个输入 \code{a/b} 相加，sum 表示本位结果，carry 表示进位。半加器没有低位进位：

\begin{lstlisting}
sum   = a XOR b
carry = a AND b
\end{lstlisting}

全加器增加 \code{carry\_in}：

\begin{lstlisting}
sum       = a XOR b XOR carry_in
carry_out = (a AND b) OR (carry_in AND (a XOR b))
\end{lstlisting}

把 64 个全加器串起来，就是 64 位 ripple-carry adder。第 0 位产生的进位传给第 1 位，第 1 位再传给第 2 位，一直传到第 63 位。

\begin{lstlisting}
carry[0] = carry_in
for i in 0..63:
  sum[i], carry[i + 1] = full_adder(a[i], b[i], carry[i])
\end{lstlisting}

这个设计直观，但关键路径很长。最高位结果必须等低位进位一路传播。若每一级全加器进位延迟是 \code{t\_fa}，64 位串行进位路径近似是 \code{64 * t\_fa}。频率高的 CPU 不能只靠这种最朴素结构。

\subsection{进位生成、传播和更快加法}

把每一位拆成 generate 和 propagate，可以更清楚地分析进位：

\begin{lstlisting}
g[i] = a[i] AND b[i]      // this bit generates carry
p[i] = a[i] XOR b[i]      // this bit propagates carry_in
carry[i + 1] = g[i] OR (p[i] AND carry[i])
\end{lstlisting}

若某一位 \code{g[i] = 1}，不管低位进位是什么，本位都会产生进位。若 \code{p[i] = 1}，本位会把低位进位传上去。carry-lookahead、carry-select、prefix adder 等结构的目标，都是减少“进位逐位等”的深度，用更多门面积换更短关键路径。

\begin{longtable}{p{0.22\linewidth}p{0.30\linewidth}p{0.36\linewidth}}
\toprule
加法器 & 关键路径特点 & 取舍 \\
\midrule
ripple-carry & 进位逐位传播，深度随位宽线性增长 & 面积小，慢 \\
carry-lookahead & 分组提前计算进位 & 面积更大，速度更快 \\
prefix adder & 树形合并 generate/propagate & 适合高性能宽位加法，布线复杂 \\
\bottomrule
\end{longtable}

OS 为什么要知道这些？不是为了手画加法器版图，而是理解“硬件速度来自结构取舍”。同样的思想会出现在页表遍历、cache 查找、队列调度里：线性等待简单但慢，提前计算和并行结构更快但更复杂。

\subsection{减法、比较和 RFLAGS}

减法可以用补码加法实现：

\begin{lstlisting}
a - b = a + bitwise_not(b) + 1
\end{lstlisting}

比较也常复用减法路径。执行 \asmcode{cmp rax, rbx} 时，硬件计算类似 \code{rax - rbx} 的结果，但不写回目的寄存器，只更新标志位。x86-64 常见条件分支会读取这些标志，例如 zero flag、sign flag、carry flag、overflow flag。

\begin{longtable}{p{0.16\linewidth}p{0.34\linewidth}p{0.38\linewidth}}
\toprule
标志 & 来源 & 软件影响 \\
\midrule
ZF & 结果是否为 0 & \asmcode{je}、循环结束、比较相等 \\
SF & 结果最高位 & 有符号比较辅助 \\
CF & 无符号进位或借位 & 无符号比较、宽整数运算 \\
OF & 有符号溢出 & 有符号比较和溢出检测 \\
\bottomrule
\end{longtable}

RFLAGS 是 CPU 架构状态的一部分。中断、异常和系统调用返回时，如果错误恢复了标志位，用户程序的后续条件分支可能走错路径。这样看，RFLAGS 不是“汇编课小细节”，而是上下文保存必须精确处理的硬件状态。

\section{ALU：把多个组合电路收束成执行单元}

ALU 可以看成一组组合电路加一个结果选择器。加法、减法、按位逻辑、移位、比较都并行或分组计算，控制信号选择最终结果和标志位。

\begin{lstlisting}
alu(a, b, op):
  add_result = a + b
  sub_result = a + (~b) + 1
  and_result = a AND b
  or_result  = a OR b
  xor_result = a XOR b
  shl_result = barrel_shift_left(a, b[0..5])

  result = mux(op,
    ADD: add_result,
    SUB: sub_result,
    AND: and_result,
    OR:  or_result,
    XOR: xor_result,
    SHL: shl_result)
  flags = compute_flags(result, op)
\end{lstlisting}

\subsection{移位器和桶形移位}

移位可以一位一位移，也可以用桶形移位器一次完成多位选择。64 位桶形移位器通常按 1、2、4、8、16、32 这些层级组织，每一层根据移位量的一位决定是否移动。它本质上是多层 MUX。

\begin{lstlisting}
stage0 = mux(shift[0], value << 1, value)
stage1 = mux(shift[1], stage0 << 2, stage0)
stage2 = mux(shift[2], stage1 << 4, stage1)
...
\end{lstlisting}

这说明高性能硬件经常用“多层选择网络”换速度。代价是面积、功耗和布线。软件里一条 \asmcode{shl rax, cl} 看起来简单，硬件里可能是一组宽 MUX 网络。

\subsection{乘法和除法为什么常有更长延迟}

加法是 ALU 的基础路径。乘法可以理解为部分积生成和求和，除法常需要迭代或复杂预测。高性能 CPU 会给乘法器专门硬件，但乘法延迟通常仍比简单整数加法长；除法更慢。

\begin{longtable}{p{0.18\linewidth}p{0.34\linewidth}p{0.36\linewidth}}
\toprule
操作 & 硬件特点 & 性能含义 \\
\midrule
ADD/SUB & 宽加法器，常见热路径 & 延迟低，吞吐高 \\
SHIFT & 多层选择网络 & 延迟中等，依赖移位器设计 \\
MUL & 部分积生成和压缩树 & 延迟更长，执行端口更少 \\
DIV & 迭代或复杂算法 & 延迟长，吞吐低 \\
\bottomrule
\end{longtable}

OS 热路径大量使用加减、位运算和比较，而谨慎使用除法。调度器、内存分配器、网络路径常把除法换成移位、乘法近似或预计算，不是写法偏好，而是硬件执行成本不同。

\section{触发器和寄存器：让机器拥有状态}

组合逻辑没有记忆。要让机器知道“当前 RIP 是多少”“当前任务是谁”“上一周期请求是否已经发出”，必须有时序逻辑。D 触发器在时钟边沿采样输入 \code{D}，把它保存为输出 \code{Q}。多个 D 触发器并排就是寄存器。

\begin{lstlisting}
on rising_clock_edge:
  if reset:
    q = 0
  else if enable:
    q = d
\end{lstlisting}

寄存器引入三个重要时序参数：

\begin{longtable}{p{0.22\linewidth}p{0.34\linewidth}p{0.32\linewidth}}
\toprule
参数 & 含义 & 违反后果 \\
\midrule
clock-to-Q & 时钟边沿后输出稳定所需时间 & 后级组合逻辑起点变晚 \\
setup time & 时钟边沿前输入必须稳定多久 & 采样错误或亚稳态 \\
hold time & 时钟边沿后输入还要保持多久 & 采样到错误变化 \\
\bottomrule
\end{longtable}

时钟周期必须满足：

\[
T_{clk} \ge T_{clk\_to\_q} + T_{comb\_max} + T_{setup} + T_{skew} + T_{margin}
\]

这就是“关键路径”概念。若一个阶段中寄存器文件读取、ALU、cache tag 比较、数据选择和写回全塞在一个周期里，\code{T\_comb\_max} 会很长，频率就上不去。流水线的本质，是在长组合路径中间插入寄存器，把一个长周期拆成多个短周期。

\section{寄存器文件：CPU 的小而快状态库}

通用寄存器不是内存数组的简单缩小版。寄存器文件需要同时支持多个读端口和写端口。例如一条整数二操作数指令常读两个源、写一个目的。系统调用 ABI 又规定 \code{rax/rdi/rsi/rdx/r10/r8/r9} 等寄存器承担系统调用号和参数。

\begin{lstlisting}
read_port0 selects source register A
read_port1 selects source register B
write_port selects destination register
write_enable controls whether result is committed
\end{lstlisting}

端口越多，寄存器文件越难做快。每个读端口都要能从所有寄存器中选择一个值，每个写端口都要把数据写到指定寄存器。乱序 CPU 还会加入寄存器重命名，把 x86-64 架构寄存器映射到更多物理寄存器，减少假依赖。

\begin{longtable}{p{0.22\linewidth}p{0.34\linewidth}p{0.32\linewidth}}
\toprule
层次 & 作用 & OS 看到什么 \\
\midrule
架构寄存器 & ISA 规定的软件可见寄存器 & trap frame、上下文切换、ABI \\
物理寄存器 & 微架构内部重命名资源 & OS 通常看不到，不逐个保存 \\
保存区 & 内存中的寄存器快照 & signal frame、switch context、pt\_regs \\
\bottomrule
\end{longtable}

上下文切换保存的是架构可见状态，而不是所有物理寄存器。中断入口保存通用寄存器，是为了让被打断代码返回后仍看到相同架构状态。理解这一点，就能区分“CPU 内部有很多临时状态”和“OS 必须保存哪些状态”。

\section{存储阵列：从寄存器到 cache 和 DRAM}

寄存器很快但面积贵，不能拿来做主内存。SRAM 用更多晶体管保存一个 bit，速度快，适合 cache、TLB 和小型片上队列。DRAM 用电容保存 bit，密度高但需要刷新，适合主内存。NAND Flash、NVMe SSD 等则提供持久存储，但访问模型完全不同。

\begin{longtable}{p{0.16\linewidth}p{0.32\linewidth}p{0.40\linewidth}}
\toprule
存储 & 硬件特点 & OS 影响 \\
\midrule
寄存器 & 最快，数量少，CPU 私有 & 上下文切换必须保存可见部分 \\
SRAM & 快，面积大，片上使用 & cache、TLB、队列、预测器 \\
DRAM & 容量大，易失，要刷新 & 物理页、NUMA、page cache、DMA buffer \\
NAND/SSD & 持久，块擦写，控制器复杂 & 文件系统、flush、FUA、磨损和错误恢复 \\
\bottomrule
\end{longtable}

\subsection{cache 作为 SRAM 结构}

cache 通常由 tag、valid、dirty、data 组成。CPU 用地址的一部分选择 set，再比较 tag，命中后读出 cache line。cache line 常见为 64 字节量级。cache 不是逐字节取一个小值，而是以 line 为单位搬运。

\begin{lstlisting}
cache_lookup(address):
  offset = address bits selecting byte in line
  index  = address bits selecting set
  tag    = remaining high bits
  for way in set[index]:
    if way.valid and way.tag == tag:
      return way.data[offset]
  miss
\end{lstlisting}

OS 设计里反复出现 cache line：自旋锁竞争会让同一 line 在 CPU 间迁移，伪共享会让不同变量互相拖慢，每 CPU 数据结构减少共享写，DMA buffer 需要考虑 CPU cache 和设备可见性。cache 是硬件，但它塑造了内核数据结构。

\subsection{TLB 作为翻译缓存}

TLB 缓存虚拟页到物理页和权限。它不是普通 data cache，因为它的内容来自页表遍历，并且带有权限、PCID、global 等语义。x86-64 地址翻译热路径大致是：

\begin{lstlisting}
virtual address
  -> canonical check
  -> TLB lookup by virtual page and PCID
  -> if miss, page walk from CR3
  -> PTE permission check
  -> physical address
\end{lstlisting}

页表项改了，内存中的 PTE 变了，但 TLB 里的旧翻译不会自动消失。于是 OS 必须执行 invalidate 和 shootdown。硬件为了快复制了状态，OS 就要负责让这些副本在关键时刻一致。

\section{数据通路：把寄存器、ALU、内存接口连起来}

有了寄存器文件、ALU、MUX、内存接口和 RIP，就可以画一个最小数据通路。以 \asmcode{add rax, rbx} 为例：

\begin{lstlisting}
1. instruction decoder selects rax as dst/src0 and rbx as src1
2. register file reads rax and rbx
3. ALU op is set to ADD
4. ALU computes sum
5. writeback MUX selects ALU result
6. register file writes rax on clock edge
7. next RIP selects sequential address
\end{lstlisting}

以 \asmcode{mov rax, qword ptr [rbx + 16]} 为例：

\begin{lstlisting}
1. decoder recognizes load with base rbx and displacement 16
2. register file reads rbx
3. address generation unit computes rbx + 16
4. D-TLB translates virtual address
5. cache or memory returns 8 bytes
6. writeback MUX selects load data
7. register file writes rax
\end{lstlisting}

这两条指令共享很多硬件：寄存器文件、ALU 或地址生成器、写回 MUX、RIP 更新路径。load 额外使用 TLB、cache、内存顺序和异常路径。OS 后面讲 page fault，就是第 4 步无法完成时，硬件把控制权交给内核。

\section{控制器：谁来产生控制信号}

数据通路回答“数据从哪里流到哪里”，控制器回答“这一拍打开哪些开关”。控制信号包括读哪个寄存器、ALU 做什么、是否访存、是否写回、下一 RIP 选哪路、当前指令是否允许在 Ring 3 执行。

\begin{longtable}{p{0.22\linewidth}p{0.34\linewidth}p{0.32\linewidth}}
\toprule
控制信号 & 控制对象 & 错误后果 \\
\midrule
reg\_read\_sel & 寄存器文件读端口 & 读错参数或地址 \\
alu\_op & ALU 功能选择 & 算错结果或标志位 \\
mem\_read/write & cache/MMU 访问 & 错读、错写或破坏内存 \\
writeback\_enable & 寄存器写回 & 状态未更新或误更新 \\
next\_rip\_select & RIP 来源 MUX & 跳错路径或无法进入异常入口 \\
privilege\_check & 特权和权限检查 & 用户态执行敏感操作 \\
\bottomrule
\end{longtable}

\subsection{硬布线控制和微码}

简单指令集可以用硬布线控制器：opcode 经过译码器直接产生控制信号。x86-64 为了长期兼容和丰富指令，会把复杂指令译成内部微操作，某些复杂或少见路径由微码协助。无论内部实现如何，软件看到的 ISA 契约必须稳定。

\begin{lstlisting}
x86-64 front-end sketch:
  fetch instruction bytes
  find instruction boundaries
  decode prefixes, opcode, ModR/M, SIB, displacement, immediate
  produce one or more micro-operations
  send uops to scheduler and execution units
\end{lstlisting}

这解释了为什么 x86-64 解码比教学固定长度指令难：一条指令可能 1 字节，也可能十几字节；操作数可能在寄存器、内存、立即数中；寻址模式可能包含 base、index、scale、displacement。OS 不写解码器，但调试、反汇编、异常 RIP、kprobes、ftrace 都依赖“指令边界可确定”。

\section{单周期 CPU：最容易理解，也最浪费}

单周期 CPU 让每条指令在一个长周期内完成所有工作。对 \asmcode{add rax, rbx}，周期要覆盖取指、译码、读寄存器、ALU、写回。对 load，周期还要覆盖地址生成、TLB/cache 访问和写回。为了让最慢指令也能完成，所有指令都被迫使用同样长的周期。

\begin{lstlisting}
T_single >= fetch + decode + reg_read + alu + tlb + cache + writeback + setup
\end{lstlisting}

假设粗略延迟如下：

\begin{longtable}{p{0.28\linewidth}p{0.24\linewidth}p{0.36\linewidth}}
\toprule
模块 & 示例延迟 & 说明 \\
\midrule
取指 I-cache & 1.0 ns & 命中时仍要 tag/data 路径 \\
译码 & 0.8 ns & x86-64 前端更复杂 \\
寄存器读 & 0.4 ns & 多端口增加负担 \\
ALU/地址生成 & 0.5 ns & 加法或地址计算 \\
D-TLB/cache & 1.2 ns & load/store 热路径 \\
写回和 setup & 0.3 ns & 写回寄存器前稳定 \\
\bottomrule
\end{longtable}

单周期 load 可能需要 4 ns 以上，于是频率上限不到 250 MHz；而简单 add 其实不需要等完整 load 路径。数字只是示意，结论重要：最长路径决定所有指令周期，硬件利用率差。

\section{多周期 CPU：把一条指令拆成状态机}

多周期 CPU 把一条指令拆成多个阶段，每个阶段用较短周期完成一部分工作。不同指令可以走不同数量的状态。例如 ALU 指令不需要 MEM 阶段，store 不需要 WB 阶段。

\begin{lstlisting}
FETCH:
  IR = memory[RIP]
  RIP = RIP + instruction_length

DECODE:
  decode IR
  read source registers

EXECUTE:
  ALU computes result or address

MEMORY:
  load or store if needed

WRITEBACK:
  update destination register if needed
\end{lstlisting}

这里第一次出现“控制器是状态机”的完整形态。控制器保存当前状态，下一状态由当前状态、opcode、异常和 ready 信号决定。内核中的驱动状态机、调度状态机、文件系统恢复状态机，形式上都和这个硬件控制器相似：状态不能跳跃，失败要有收敛路径。

\section{流水线 CPU：用寄存器切开关键路径}

流水线在阶段之间插入流水线寄存器，让 IF、ID、EX、MEM、WB 同时处理不同指令。它提高吞吐，不一定降低单条指令从开始到完成的延迟。

\begin{lstlisting}
cycle 1: inst1 IF
cycle 2: inst1 ID   inst2 IF
cycle 3: inst1 EX   inst2 ID   inst3 IF
cycle 4: inst1 MEM  inst2 EX   inst3 ID   inst4 IF
cycle 5: inst1 WB   inst2 MEM  inst3 EX   inst4 ID   inst5 IF
\end{lstlisting}

流水线周期由最慢阶段决定：

\[
T_{pipe} \ge \max(T_{IF}, T_{ID}, T_{EX}, T_{MEM}, T_{WB}) + T_{pipe\_reg}
\]

流水线寄存器也有成本。切得太浅，频率上不去；切得太深，分支错误和异常冲刷成本增加，流水线寄存器本身耗电也增加。现代 x86-64 CPU 的前端、乱序后端、load/store、提交都比五级模型复杂，但关键思想不变：用状态边界把长组合路径拆短。

\subsection{数据 hazard}

相邻指令有依赖时，流水线会遇到数据 hazard：

\begin{lstlisting}
add rax, rbx
sub rcx, rax
\end{lstlisting}

第二条需要第一条的新 \code{rax}。如果等第一条写回再读，流水线要停。硬件可用 forwarding 把执行结果直接送到后续执行单元，减少等待。load-use 更难，因为 load 数据要到 cache 返回后才知道。

\begin{lstlisting}
if EX/MEM.rd == ID/EX.rs and EX/MEM.writes_register:
  operand = EX/MEM.forwarded_result
else if MEM/WB.rd == ID/EX.rs and MEM/WB.writes_register:
  operand = MEM/WB.forwarded_result
else:
  operand = register_file_value
\end{lstlisting}

\subsection{控制 hazard 和分支预测}

分支决定下一 RIP。如果等分支完全解析再取指，前端会空等。分支预测器猜方向和目标，先把预测路径送进流水线。猜错时，硬件丢弃年轻指令，从正确 RIP 重新取指。

\begin{lstlisting}
if prediction_wrong:
  squash younger pipeline entries
  restore front-end RIP to correct target
  restart fetch
\end{lstlisting}

分支预测是性能优化，也带来安全问题。被冲刷的预测路径不会提交架构结果，但可能改变 cache 和预测器状态。OS 的隔离和安全补丁要考虑这种微架构痕迹。

\subsection{结构 hazard}

结构 hazard 是两个阶段争同一个硬件资源。例如取指和数据访问若共用同一个单端口内存，会冲突；多个执行单元不足时，后端会排队。硬件通过分离 I-cache/D-cache、多端口结构、队列和调度器缓解。

这对 OS 性能分析很重要。程序慢不一定是算法指令数多，也可能是前端取指不足、load/store 单元堵塞、执行端口竞争、cache miss 或 TLB miss。perf 看到 IPC 低时，要能把它落到这些硬件资源上。

\section{乱序执行：保持 ISA 顺序，内部尽量并行}

乱序执行把“何时执行”和“何时提交”分开。指令进入 reorder buffer，源操作数准备好就可以发射到执行单元；结果先写到临时位置；最终按程序顺序提交到架构状态。

\begin{lstlisting}
front-end decodes instructions into uops
rename maps architectural registers to physical registers
scheduler wakes uops whose operands are ready
execution units run ready uops out of order
ROB commits completed uops in program order
\end{lstlisting}

这样可以在 cache miss 等慢事件发生时继续执行无关工作。但 OS 依赖的是提交边界：异常必须精确，系统调用必须在受控边界进入，调试单步必须看到可解释 RIP，上下文切换保存的是已经提交或可恢复的架构状态。

\begin{longtable}{p{0.22\linewidth}p{0.34\linewidth}p{0.32\linewidth}}
\toprule
内部结构 & 作用 & OS 依赖的外部契约 \\
\midrule
rename table & 消除假依赖，提高并行度 & 架构寄存器语义不变 \\
scheduler & 选择可执行 uop & 指令结果最终按 ISA 呈现 \\
ROB & 记录顺序和异常 & precise exception、可重试 fault \\
store buffer & 缓冲写入 & 需要内存屏障控制可见顺序 \\
\bottomrule
\end{longtable}

\section{异常和中断硬件：CPU 怎样受控交还给 OS}

异常和中断不是在 CPU 外面“调用内核函数”。硬件必须在可恢复边界停止当前流，保存返回所需状态，查 IDT，必要时切换到 TSS 指定的内核栈，然后跳到 handler。

\begin{lstlisting}
event_entry(vector):
  finish or roll back to precise boundary
  determine whether privilege changes
  if Ring 3 -> Ring 0:
    switch RSP to TSS.RSP0 or IST stack
  push return RIP, CS, RFLAGS, RSP, SS as required
  push error code for selected exceptions
  RIP = IDT[vector].handler
\end{lstlisting}

这条路径依赖前面所有硬件：RIP 是寄存器，IDT 查找是表项读取和权限检查，TSS 提供栈，RFLAGS 记录状态，流水线和 ROB 提供精确边界。OS 的 trap frame 不是软件随便造的结构，而是硬件入口事实和汇编保存约定合成的证据。

\section{MMU 和页表硬件：地址也要通过电路决策}

虚拟地址不是直接送到 DRAM。x86-64 CPU 要先检查 canonical 形式，再查 TLB，miss 时从 CR3 指向的顶级页表开始 page walk。每级页表项都要检查 present、权限、地址位和页大小等字段。

\begin{lstlisting}
translate(va, access, cpl):
  require canonical(va)
  if TLB has (pcid, virtual_page):
    check cached permissions
    return physical_page + offset

  pml4e = read_phys(CR3.base + pml4_index(va) * 8)
  pdpte = read_phys(pml4e.base + pdpt_index(va) * 8)
  pde   = read_phys(pdpte.base + pd_index(va) * 8)
  pte   = read_phys(pde.base + pt_index(va) * 8)
  check permission bits
  fill TLB
  return pte.page + offset
\end{lstlisting}

page walk 本身也是内存访问，可能命中 cache，也可能很慢。因此 TLB 是性能关键。修改页表后必须失效旧 TLB 项，是因为硬件为了快保存了翻译副本。COW、\code{mprotect}、\code{munmap}、进程退出都要处理这个副本生命周期。

\section{总线和 ready/valid：模块之间怎样说话}

CPU 内部和外部模块不能只靠“函数返回”。硬件常用 ready/valid 协议表达一次传输：发送方在数据有效时拉高 valid，接收方能接受时拉高 ready，二者同时为 1 的周期完成一次握手。

\begin{lstlisting}
transfer happens when valid && ready
sender must keep data stable while valid && !ready
receiver may apply backpressure by lowering ready
\end{lstlisting}

这个小协议解释了很多 OS 现象。设备队列满时就是硬件或驱动层 backpressure；block layer 要限制 queue depth；网络栈要处理发送队列和接收 ring；io\_uring 要区分 submission queue 和 completion queue。硬件接口天然就是异步状态机，不是同步函数调用。

\section{从硬件模块到 OS 抽象}

把本章的硬件模块映射回 OS，可以看到每个抽象都有电路根。

\begin{longtable}{p{0.24\linewidth}p{0.34\linewidth}p{0.30\linewidth}}
\toprule
硬件层 & 直接能力 & OS 抽象或约束 \\
\midrule
逻辑门和 MUX & 条件选择、控制信号 & 特权入口只能走受控路径 \\
加法器和 ALU & 地址计算、比较、标志位 & 指针、页内偏移、RFLAGS、条件分支 \\
触发器和寄存器 & 保存当前状态 & trap frame、context switch、signal frame \\
寄存器文件 & 快速操作数和 ABI 位置 & syscall 参数、返回值、调用约定 \\
cache/TLB SRAM & 缓存数据和地址翻译 & 局部性、TLB shootdown、cache line 竞争 \\
流水线和 ROB & 高吞吐和精确提交 & precise fault、可重试 page fault \\
MMU/page walker & 地址翻译和权限检查 & 地址空间、COW、mmap、隔离 \\
IDT/TSS/APIC & 异常、中断、栈切换 & syscall、timer、IRQ、IPI、调度 \\
ready/valid/队列 & 异步通信和背压 & DMA ring、completion、等待队列 \\
\bottomrule
\end{longtable}

因此，后面讲 OS 时要持续追问：这个抽象最后落在哪个硬件状态上？谁能改它？什么时候采样？旧副本在哪里？失败时哪个入口拿到证据？能回答这些问题，才算把硬件和 OS 真正合起来学。

\section{硬件章节验收}

\begin{rubricbox}
\begin{itemize}
\tightlist
\item 能从 NAND/NOR 说明为什么逻辑门能组合成任意布尔函数。
\item 能画出半加器、全加器，并解释 64 位 ripple-carry 的关键路径。
\item 能说明 carry-lookahead 为什么用更多硬件换更短时延。
\item 能把 \asmcode{cmp rax, rbx} 解释为减法路径和 RFLAGS 更新。
\item 能说明寄存器的 setup/hold/clock-to-Q 如何限制时钟周期。
\item 能把寄存器文件、ALU、MUX、RIP 更新连成一条 \asmcode{add rax, rbx} 数据通路。
\item 能解释单周期、多周期、流水线 CPU 的时延取舍。
\item 能说明 forwarding、stall、flush、ROB 和 precise exception 分别解决什么问题。
\item 能解释 TLB 为什么不是页表本体，以及页表修改后为什么要 invalidate。
\item 能把 ready/valid、队列和 backpressure 映射到 DMA ring、completion 和等待队列。
\end{itemize}
\end{rubricbox}
